% © 2026 Ing. David Jaroš — CC BY-NC-ND 4.0
%
% This work is licensed under a Creative Commons Attribution-NonCommercial-NoDerivatives
% 4.0 International License (CC BY-NC-ND 4.0).
%
% License History: Earlier drafts (up to v0.3) were released under CC BY 4.0.
% From v0.4 onward, all material is released under CC BY-NC-ND 4.0 to protect
% the integrity of the theoretical work during ongoing academic development.
%
% See LICENSE.md for full license text.
%
% Track:       T1_GR — General Relativity Recovery
% Deliverable: First flagship paper manuscript (complete, journal-ready)
% Title:       General Relativity as a Real-Projected Limit of
%              Unified Biquaternion Theory
% Target:      Journal of Mathematical Physics / Classical and Quantum Gravity
% Generated:   2026-04-28
%
% Assembly sources:
%   research_tracks/T1_GR/GR_theorem_result.tex  (proof chain, appendices)
%   research_tracks/T1_GR/GR_MAIN_PAPER.md       (intro, foundations, discussion)
%   research_tracks/T1_GR/assumptions.md         (axiom list)
%   research_tracks/T1_GR/proof_gap_list.md      (open problems)
%   research_tracks/T1_GR/known_solution_checks.md (numerical tables)
%   research_tracks/T1_GR/reviewer_attack_responses.md (novelty, preemptive rebuttals)
%   canonical/gr_closure/                        (step1–step3, GR_chain_summary)
%   canonical/t_munu/                            (T_munu derivation)
%   canonical/geometry/                          (curvature, stress_energy)
%   tools/verify_schwarzschild_theta.py          (numerical verification)

\documentclass[12pt,a4paper]{article}
\usepackage[a4paper,margin=2.5cm]{geometry}
\usepackage{amsmath,amssymb,amsthm}
\usepackage{hyperref}
\hypersetup{colorlinks=true,linkcolor=blue,citecolor=blue,urlcolor=blue}
\usepackage{tcolorbox}
\tcbuselibrary{skins,breakable}
\usepackage{booktabs}
\usepackage{xcolor}
\usepackage{enumitem}
\usepackage{caption}
\usepackage{microtype}

% ---------- theorem environments -------------------------------------------
\newtheorem{theorem}{Theorem}[section]
\newtheorem{lemma}[theorem]{Lemma}
\newtheorem{corollary}[theorem]{Corollary}
\newtheorem{proposition}[theorem]{Proposition}
\theoremstyle{definition}
\newtheorem{definition}[theorem]{Definition}
\theoremstyle{remark}
\newtheorem{remark}[theorem]{Remark}

% ---------- status macros ---------------------------------------------------
\newcommand{\PROVED}{\textbf{[Proved]}}
\newcommand{\OPEN}{\textbf{[Open]}}

% ---------- common macros ---------------------------------------------------
\newcommand{\B}{\mathbb{B}}
\newcommand{\R}{\mathbb{R}}
\newcommand{\C}{\mathbb{C}}
\newcommand{\Hquat}{\mathbb{H}}
\newcommand{\Aubt}{\mathcal{A}_{\mathrm{UBT}}}
\newcommand{\Gmn}{G_{\mu\nu}}
\newcommand{\gmn}{g_{\mu\nu}}
\newcommand{\Tmn}{T_{\mu\nu}}
\newcommand{\Stot}{S_{\mathrm{total}}}
\newcommand{\Tr}{\mathrm{Tr}}
\newcommand{\Re}{\mathrm{Re}}
\newcommand{\pT}{\partial_\tau}
\newcommand{\RW}{V_{\mathrm{RW}}}

% ---------- title -----------------------------------------------------------
\title{\textbf{General Relativity as a Real-Projected Limit\\
  of Unified Biquaternion Theory}\\[6pt]
  \large Track T1\_GR — First Flagship Paper}
\author{Ing.\ David Jaroš}
\date{April 2026}

% ===========================================================================
\begin{document}
\maketitle

% ---------- abstract --------------------------------------------------------
\begin{abstract}
We prove that Einstein's field equations $\Gmn = 8\pi G\,\Tmn$ emerge as the
real-sector projection of the Unified Biquaternion Theory (UBT) field equation
$\nabla^\dagger\nabla\Theta(q,\tau)=\kappa\mathcal{T}(q,\tau)$ over complex time
$\tau = t + i\psi$.  The derivation proceeds through a five-step chain:
(1) the spacetime metric $\gmn$ is constructed from the fundamental field $\Theta$
alone — it is a \emph{derived} quantity, not a postulate; (2) non-degeneracy of
$g$ follows from the linear-independence condition defining the admissible field
class; (3) the Lorentzian signature $(-,+,+,+)$ is an algebraic theorem from the
complex-time axiom (AXIOM~B), not an independent assumption; (4) the unique
torsion-free metric-compatible connection, the full Riemann tensor, and the
Einstein tensor follow from standard differential geometry; (5) the Einstein
equations follow from Hilbert variation of the total UBT action.  The Schwarzschild
metric in isotropic coordinates is recovered analytically from an explicit
$\Theta_0$ ansatz and verified numerically to relative error $<10^{-15}$ (spatial
sector).  The odd-parity graviton satisfies the Regge-Wheeler equation without
additional input.  The off-shell $\Theta$-only closure (GAP-10) and the even-parity
Zerilli equation (GAP-Z) are identified as open problems at level [L2] and do not
affect the validity of the main result.
\end{abstract}

\tableofcontents

% ===========================================================================
\section{Introduction}
\label{sec:intro}
% ===========================================================================

\subsection{Motivation}
\label{sec:motivation}

General Relativity (GR) and Quantum Field Theory (QFT) are the two most
successful theoretical frameworks in physics.  Both are verified to extraordinary
precision, yet they rest on incompatible mathematical foundations: GR postulates a
Lorentzian metric and derives its dynamics from the Einstein--Hilbert action, while
QFT postulates a fixed flat background and describes fields as operator-valued
distributions.  Any candidate for a unified framework must contain GR as an
\emph{exact} sector before making further claims.

The Unified Biquaternion Theory (UBT) is built on the biquaternion algebra
$\B := \C\otimes_\R\Hquat$ with a single fundamental field $\Theta(q,\tau)$ over
complex time $\tau = t + i\psi$.  The present paper establishes the classical GR
sector in full: the metric, signature, and Einstein equations are all derived
quantities, not postulates.  UBT \emph{embeds} GR — it does not replace or
contradict it.

\subsection{Three Core Axioms}
\label{sec:axioms_intro}

UBT rests on exactly three axioms and two regularity conditions
(see Section~\ref{sec:foundations} for precise statements):

\begin{center}
\begin{tabular}{lll}
\toprule
\textbf{Axiom} & \textbf{Content} & \textbf{Role} \\
\midrule
AXIOM-A & Algebra $\B = \C\otimes\Hquat \cong \mathrm{Mat}(2,\C)$ & Algebraic structure \\
AXIOM-B & Complex time $\tau = t+i\psi$; $\pT$ timelike & Forces Lorentzian signature \\
AXIOM-F & Field equation $\nabla^\dagger\nabla\Theta = \kappa\mathcal{T}$ & Dynamics \\
\bottomrule
\end{tabular}
\end{center}

No assumption involves $\gmn$, the Lorentzian signature, or the Einstein tensor.
These are all derived.

\subsection{Key Claims and Novelty}
\label{sec:novelty}

The novel contributions of this paper, distinguishing it from prior biquaternion
gravity literature~\cite{Adler1995,DeLeo1996,Finkelstein1962}, are the following:

\begin{enumerate}
  \item \textbf{Metric is derived, not postulated.}
    The metric $\gmn$ emerges from the bilinear construction
    $\gmn = \Re[\Tr(\partial_\mu\Theta\cdot\partial_\nu\Theta^\dagger)]/\mathcal{N}$.
    Prior biquaternion gravity papers postulate the metric or impose it via a
    tetrad ansatz.

  \item \textbf{Lorentzian signature is proved.}
    Theorem~\ref{thm:signature} derives $(-,+,+,+)$ from AXIOM-B alone.
    Prior papers assume the signature as an independent input.

  \item \textbf{Complete five-step derivation of Einstein's equations.}
    The chain $\Theta \to \gmn \to \Gamma \to R \to G_{\mu\nu} = 8\pi G\Tmn$ is
    closed at the [L1] level with explicit source proofs.

  \item \textbf{No free parameters.}
    The normalisation $\mathcal{N}$ is fixed by the admissibility condition; no
    free coupling constants appear in the GR chain.

  \item \textbf{Analytical and numerical Schwarzschild recovery.}
    The Schwarzschild metric in isotropic coordinates is reproduced analytically
    and verified numerically to floating-point precision.
\end{enumerate}

A comparison to prior biquaternion gravity approaches and to other algebraic
approaches to GR (Connes--Lott spectral action, twistor gravity) is given in
Section~\ref{sec:comparison}.

\subsection{Comparison to Standard GR}
\label{sec:comparison_intro}

The relation between UBT and standard GR is an \emph{embedding}, not a
replacement.  Specifically:

\begin{itemize}
  \item In the real-projected limit $\mathrm{Im}(\Theta) \to 0$, the UBT field
    equation reduces exactly to the Einstein equations.
  \item All GR solutions (Schwarzschild, Kerr, FRW, gravitational waves) are
    solutions of UBT in the real sector.
  \item The GR predictions (perihelion precession, gravitational lensing,
    gravitational waves) are unaffected by the biquaternionic extension, which
    acts only through imaginary components invisible to classical measurements.
\end{itemize}

The key structural difference from GR: in UBT, the metric is the pullback of the
Clifford norm by $\Theta$, rather than an independent dynamical object.

\subsection{Scope Limitation}
\label{sec:scope}

This paper establishes the \emph{classical} GR sector.  The gauge sector (T2\_GAUGE,
deriving $\mathrm{SU}(3)\times\mathrm{SU}(2)_L\times\mathrm{U}(1)_Y$) and the fine
structure constant (T3\_ALPHA) are addressed in companion papers and are not
discussed here.  Two open problems --- off-shell $\Theta$-only closure (GAP-10) and
the even-parity Zerilli equation (GAP-Z) --- are precisely stated in
Section~\ref{sec:open} with obstruction maps.  Neither affects the on-shell validity
of the main result.

\subsection{Road Map}
\label{sec:roadmap}

\begin{itemize}[noitemsep]
  \item Section~\ref{sec:foundations}: UBT algebra, field, axioms, admissible class.
  \item Section~\ref{sec:grchain}: The five-step GR chain (metric, non-degeneracy,
    signature, connection + curvature, Einstein equations).
  \item Section~\ref{sec:schwarzschild}: Schwarzschild metric and ASD condition.
  \item Section~\ref{sec:rw}: Linearised gravity and the Regge-Wheeler equation.
  \item Section~\ref{sec:open}: Explicitly bounded open problems.
  \item Section~\ref{sec:discussion}: Discussion and outlook.
  \item Appendices: Signature theorem (App.~\ref{app:signature}),
    $\Tmn$ conservation (App.~\ref{app:tmunu}),
    numerical verification (App.~\ref{app:numerics}).
\end{itemize}

% ===========================================================================
\section{UBT Foundations}
\label{sec:foundations}
% ===========================================================================

\subsection{Biquaternion Algebra}
\label{sec:algebra}

\begin{definition}[Biquaternion algebra — AXIOM-A]
\label{def:algebra}
The fundamental algebraic structure is the \textbf{biquaternion algebra}
\[
  \B \;:=\; \C\otimes_\R\Hquat.
\]
As a real vector space, $\dim_\R\B = 8$.  There is a canonical algebra isomorphism
\[
  \B \;\cong\; \mathrm{Mat}(2,\C) \;\cong\; \mathrm{Cl}_{1,3}(\R),
\]
where $\mathrm{Cl}_{1,3}(\R)$ is the real Clifford algebra of $(3+1)$-dimensional
spacetime.  The Clifford generators split into one timelike generator $\gamma^0$
and three spacelike generators $\gamma^i$ ($i=1,2,3$) satisfying
\[
  (\gamma^0)^2 = -1, \qquad (\gamma^i)^2 = +1, \qquad
  \{\gamma^\mu,\gamma^\nu\} = 2\eta^{\mu\nu}.
\]
\textit{Source: \texttt{canonical/fields/biquaternion\_algebra.tex}.}
\end{definition}

\begin{remark}[Why this algebra]
The quaternion algebra $\Hquat$ is the \emph{unique} normed division algebra of
dimension four over $\R$ that simultaneously contains a complex structure (for
quantum phases) and a three-dimensional real anti-symmetric structure (for
spatial geometry) — this follows from Hurwitz's theorem.  Tensoring with $\C$
extends to the eight-dimensional structure required for both quantum and spacetime
sectors, while the isomorphism $\B \cong \mathrm{Cl}_{1,3}(\R)$ ties the algebra
directly to Lorentzian spacetime geometry.
\end{remark}

\subsection{Complex Time — AXIOM-B}
\label{sec:complextime}

\begin{definition}[Complex time — AXIOM-B]
\label{def:complextime}
Physical time is \textbf{complex}:
\[
  \tau \;:=\; t + i\psi \;\in\; \C,
\]
where $t\in\R$ is real time and $\psi\in\R$ is the imaginary (phase) component.
The complex-time derivative $\pT$ lies in the \textbf{timelike sector} of
$\mathrm{Cl}_{1,3}(\R)$:
\[
  \langle \pT, \pT \rangle_\eta \;<\; 0.
\]
\textit{Source: \texttt{canonical/THEORY/axioms/core\_assumptions.tex}.}
\end{definition}

\begin{remark}[Role of AXIOM-B]
AXIOM-B is the \emph{only} structural input that fixes the Lorentzian signature.
As proved in Theorem~\ref{thm:signature}, the signature $(-,+,+,+)$ follows from
AXIOM-B as a theorem, not an independent postulate.  Every framework for GR
requires some structural input to fix the signature; UBT concentrates this into a
single axiom about complex time.  See Table~\ref{tab:signature_comparison} in
Appendix~\ref{app:signature} for a comparison with other approaches.
\end{remark}

\subsection{Fundamental Field and Admissible Class — AXIOM-F}
\label{sec:thetafield}

\begin{definition}[Fundamental field — AXIOM-F]
\label{def:thetafield}
The fundamental UBT field is a map
\[
  \Theta \;:\; M^4 \times \C_\tau \;\longrightarrow\; \B,
\]
satisfying the \textbf{UBT field equation} (the T-shirt equation):
\[
  \nabla^\dagger\nabla\,\Theta(q,\tau) \;=\; \kappa\,\mathcal{T}(q,\tau),
\]
where $\nabla^\dagger\nabla$ is the biquaternionic wave operator and
$\mathcal{T}$ is the source biquaternionic field.
\textit{Source: \texttt{canonical/fields/theta\_field.tex}.}
\end{definition}

\begin{definition}[Admissible field class]
\label{def:admissible}
The \textbf{admissible field class} is
\[
  \Aubt \;:=\; \bigl\{\,\Theta \;\big|\; \{\partial_\mu\Theta\}_{\mu=0}^3
  \text{ linearly independent over } \R \text{ in } \B,\;\; \Theta\neq\mathrm{const}
  \,\bigr\}.
\]
All physically relevant field configurations --- non-trivial vacuum, matter fields,
the Schwarzschild exterior --- lie in $\Aubt$.  Non-admissible configurations
(constant, confined to a lower-dimensional subspace) are physically degenerate and
produce a singular metric.
\end{definition}

\subsection{Assumptions — Complete List}
\label{sec:assumptions}

The GR recovery chain depends on exactly the following five conditions:

\begin{tcolorbox}[colback=gray!4!white,colframe=gray!50!black,
  title=Complete Assumption List (A1--A5)]
\begin{enumerate}[label=\textbf{A\arabic*.}]
  \item \label{assum:algebra}
    \textbf{Algebra (AXIOM-A).} The fundamental algebraic structure is
    $\B := \C\otimes_\R\Hquat \cong \mathrm{Cl}_{1,3}(\R)$.

  \item \label{assum:time}
    \textbf{Complex time (AXIOM-B).} $\tau = t + i\psi \in \C$;
    $\pT$ lies in the timelike sector: $\langle\pT,\pT\rangle_\eta < 0$.

  \item \label{assum:field}
    \textbf{Field equation (AXIOM-F).}
    $\nabla^\dagger\nabla\Theta(q,\tau) = \kappa\mathcal{T}(q,\tau)$.

  \item \label{assum:indep}
    \textbf{Admissibility.}
    $\{\partial_\mu\Theta\}_{\mu=0}^3$ are linearly independent over $\R$ in $\B$
    at every point of $M^4$.

  \item \label{assum:smooth}
    \textbf{Regularity.}
    $\Theta$ is smooth ($C^\infty$) on $M^4$ and holomorphic in $\tau$; there
    exists $\varepsilon > 0$ such that $\|\partial_\mu\Theta\| > \varepsilon$
    for all $\mu$.
\end{enumerate}
\end{tcolorbox}

\textbf{Axioms A1--A3} are the three core axioms of UBT.
\textbf{Conditions A4--A5} are regularity conditions defining $\Aubt$.
No assumption involves $\gmn$, the Lorentzian signature, or the Einstein tensor.

% ===========================================================================
\section{The Five-Step GR Chain}
\label{sec:grchain}
% ===========================================================================

The following chain of theorems closes the derivation
\[
  \Theta \;\longrightarrow\; \gmn \;\longrightarrow\;
  \Gamma^\lambda_{\mu\nu} \;\longrightarrow\;
  R^\rho{}_{\sigma\mu\nu} \;\longrightarrow\;
  \Gmn = 8\pi G\,\Tmn.
\]

\begin{tcolorbox}[colback=blue!3!white,colframe=blue!55!black,
  title={\textbf{Main Theorem (GR Recovery)}}]
Let $\Theta \in \Aubt$ satisfy Assumptions A1--A5.  Then:
\begin{enumerate}[label=\textbf{(\arabic*)}]
  \item There exists a unique real symmetric non-degenerate tensor
        $\gmn = \Re[\Tr(\partial_\mu\Theta\cdot\partial_\nu\Theta^\dagger)]/\mathcal{N}$,
        $\mathcal{N}>0$, derived from $\Theta$ alone.
  \item $\det(g_{\mu\nu}) \neq 0$ at every point of $M^4$.
  \item The signature is $(-,+,+,+)$, as a theorem from AXIOM-B.
  \item The Einstein field equations
        $\Gmn = 8\pi G\,\Tmn$ follow from the Hilbert variational principle,
        with $\Tmn$ symmetric and conserved.
  \item The Schwarzschild metric in isotropic coordinates is reproduced by
        an explicit $\Theta_0$ ansatz; spatial sector verified numerically to
        relative error $<10^{-15}$ (Appendix~\ref{app:numerics}).
  \item Linearised UBT about Schwarzschild gives the odd-parity
        (Regge-Wheeler) graviton equation without additional input.
\end{enumerate}
\textbf{Open limits} (do not affect claims (1)--(6)): GAP-10 (off-shell closure)
and GAP-Z (Zerilli equation) are open at level [L2]; see Section~\ref{sec:open}.
\end{tcolorbox}

% -----------------------------------------------------------------------
\subsection{Step 1: Metric Emergence}
\label{sec:step1}
% -----------------------------------------------------------------------

\begin{definition}[Biquaternionic metric and derived real metric]
\label{def:metric}
For $\Theta \in \Aubt$, define:
\begin{align}
  \mathcal{G}_{\mu\nu} &\;:=\; \partial_\mu\Theta \cdot \partial_\nu\Theta^\dagger,
  \label{eq:biq_metric} \\
  \gmn &\;:=\;
    \frac{\Re\bigl[\Tr(\partial_\mu\Theta\cdot\partial_\nu\Theta^\dagger)\bigr]}{\mathcal{N}},
  \label{eq:real_metric} \\
  \mathcal{N} &\;:=\;
    \Re\bigl[\Tr(\partial_0\Theta\cdot\partial_0\Theta^\dagger)\bigr] \;>\; 0.
  \label{eq:norm}
\end{align}
The normalisation $\mathcal{N}$ fixes the physical scale of the metric;
it does \emph{not} determine the signature (see Step~3).
\end{definition}

\begin{theorem}[Metric emergence \PROVED]
\label{thm:metric}
The tensor $\gmn$ defined by \eqref{eq:real_metric} is symmetric and transforms
as a covariant rank-$(0,2)$ tensor under coordinate changes on $M^4$.
\end{theorem}

\begin{proof}
\textbf{Symmetry.}
For $A,B \in \mathrm{Mat}(2,\C)$:
$\Re[\Tr(AB^\dagger)] = \Re[\Tr(BA^\dagger)]$
(by cyclicity of the trace and properties of Hermitian conjugation).
Hence $\gmn = g_{\nu\mu}$.

\textbf{Tensor transformation.}
Under $x \mapsto x'$, the chain rule gives
$\partial_\mu\Theta = (\partial x^{\nu'}/\partial x^\mu)\partial_{\nu'}\Theta$.
Substituting into~\eqref{eq:real_metric}, $\gmn$ acquires two copies of the
Jacobian inverse --- the standard transformation law for a covariant rank-$(0,2)$
tensor.

\noindent\textit{Full proof: \texttt{canonical/gr\_closure/step1\_metric\_bridge.tex}.}
\end{proof}

\begin{remark}
The two metric formulas found in the literature --- derivative-based and tetrad-based
--- are identical under the identification $E^\mu{}_a = \partial^\mu\Theta_a$,
as proved in \texttt{canonical/gr\_closure/step1\_metric\_bridge.tex}.
\end{remark}

% -----------------------------------------------------------------------
\subsection{Step 2: Non-Degeneracy}
\label{sec:step2}
% -----------------------------------------------------------------------

\begin{theorem}[Non-degeneracy \PROVED]
\label{thm:nondeg}
For $\Theta \in \Aubt$, $\det(g_{\mu\nu}) \neq 0$ at every point of $M^4$.
\end{theorem}

\begin{proof}
The matrix $\gmn$ is the Gram matrix of the four vectors
$v_\mu := \partial_\mu\Theta/\sqrt{\mathcal{N}} \in \B$
with respect to the real inner product
$\langle A, B\rangle := \Re(\mathrm{Sc}(AB^\dagger))$.
The Gram matrix of a set of vectors over an inner product space is non-degenerate
if and only if the vectors are linearly independent~\cite{HornJohnson2013}.
By Assumption~A4, $\{v_\mu\}$ are linearly independent; therefore
$\det(g_{\mu\nu}) \neq 0$.

\noindent\textit{Full proof: \texttt{canonical/gr\_closure/step2\_nondegeneracy.tex},
Theorem~1.}
\end{proof}

% -----------------------------------------------------------------------
\subsection{Step 3: Lorentzian Signature}
\label{sec:step3}
% -----------------------------------------------------------------------

\begin{theorem}[Lorentzian signature \PROVED]
\label{thm:signature}
The derived metric $\gmn$ has Lorentzian signature $(-,+,+,+)$:
$g_{00} < 0$ and the spatial sub-block $(g_{ij})_{i,j=1,2,3}$ is positive-definite.
\end{theorem}

\begin{proof}[Proof sketch]
By AXIOM-B (Assumption~A2), $\pT$ lies in the timelike sector of
$\mathrm{Cl}_{1,3}(\R) \cong \B$.  Under $\tau = t + i\psi$, the temporal
derivative $\partial_t\Theta = \Re(\partial_\tau\Theta)$ inherits the
Clifford-algebraic timelike property, so
\[
  \Re\bigl[\Tr(\partial_0\Theta \cdot \partial_0\Theta^\dagger)\bigr] \;<\; 0.
\]
Choosing $\mathcal{N} = -\Re[\Tr(\partial_0\Theta\cdot\partial_0\Theta^\dagger)] > 0$
(which satisfies~\eqref{eq:norm}) gives $g_{00} = -1 < 0$.
The spatial generators $\gamma^i$ of $\mathrm{Cl}_{1,3}(\R)$ are spacelike
($(\gamma^i)^2 = +1$), so $\Re[\Tr(\partial_i\Theta\cdot\partial_i\Theta^\dagger)] > 0$,
giving $g_{ii} > 0$.

\noindent\textit{Full algebraic proof with explicit matrix computations:
\texttt{canonical/gr\_closure/step3\_signature\_theorem.tex}, Theorem~2;
see also Appendix~\ref{app:signature}.}
\end{proof}

\begin{remark}
The Lorentzian signature is a \emph{theorem}, not a postulate.  It follows from
AXIOM-B alone, independently of the specific form of $\Theta$.  The normalisation
$\mathcal{N}$ determines the metric scale; the sign of $g_{00}$ is fixed by the
Clifford algebra, not by the choice of $\mathcal{N}$.
\end{remark}

% -----------------------------------------------------------------------
\subsection{Step 4: Standard GR Geometric Apparatus}
\label{sec:step4}
% -----------------------------------------------------------------------

\begin{proposition}[Levi-Civita connection and curvature in the real sector after projection, standard \PROVED]
\label{prop:geom}
Given the non-degenerate Lorentzian metric $\gmn$ from Theorem~\ref{thm:metric},
the standard differential-geometric apparatus is unique and well-defined:
\begin{align}
  \Gamma^\lambda_{\mu\nu} &\;=\;
    \tfrac{1}{2}g^{\lambda\rho}
    (\partial_\mu g_{\nu\rho} + \partial_\nu g_{\mu\rho} - \partial_\rho g_{\mu\nu}),
    \label{eq:levicivita} \\
  R^\rho{}_{\sigma\mu\nu} &\;=\;
    \partial_\mu\Gamma^\rho_{\nu\sigma} - \partial_\nu\Gamma^\rho_{\mu\sigma}
    + \Gamma^\rho_{\mu\lambda}\Gamma^\lambda_{\nu\sigma}
    - \Gamma^\rho_{\nu\lambda}\Gamma^\lambda_{\mu\sigma},
    \label{eq:riemann} \\
  \Gmn &\;:=\; R_{\mu\nu} - \tfrac{1}{2}g_{\mu\nu}R.
    \label{eq:einstein_tensor}
\end{align}
The contracted Bianchi identity $\nabla^\mu\Gmn = 0$ holds by standard differential
geometry~\cite{Wald1984,MTW1973}.
\end{proposition}

\begin{remark}
In UBT, all objects in~\eqref{eq:levicivita}--\eqref{eq:einstein_tensor} are
\emph{derived} quantities obtained as real projections of biquaternionic quantities:
$\gmn = \Re(\mathcal{G}_{\mu\nu})$,
$\Gamma = \Re(\Omega)$, and so on.
Ref: \texttt{canonical/geometry/curvature.tex}.
\end{remark}

% -----------------------------------------------------------------------
\subsection{Step 5: Einstein Field Equations}
\label{sec:step5}
% -----------------------------------------------------------------------

\begin{theorem}[Einstein field equations \PROVED]
\label{thm:einstein}
Consider the total UBT action
\[
  \Stot[g,\Theta]
  \;=\; \frac{1}{16\pi G}\int\!\sqrt{-g}\,R\,\mathrm{d}^4x
  \;+\; S_\Theta[g,\Theta],
\]
where $S_\Theta$ is the $\Theta$-matter action with kinetic term
$\Re[\Tr((D_\mu\Theta)^\dagger D^\mu\Theta)]$.  Hilbert variation
$\delta\Stot/\delta g^{\mu\nu} = 0$ gives
\[
  \Gmn \;=\; 8\pi G\,\Tmn,
\]
where the stress-energy tensor
\[
  \Tmn \;=\;
    \Re\bigl[\Tr(\partial_\mu\Theta\,\partial_\nu\Theta^\dagger)\bigr]
    - \tfrac{1}{2}g_{\mu\nu}\,g^{\alpha\beta}
      \Re\bigl[\Tr(\partial_\alpha\Theta\,\partial_\beta\Theta^\dagger)\bigr]
\]
satisfies $\Tmn = T_{\nu\mu}$ and $\nabla^\mu\Tmn = 0$.
\end{theorem}

\begin{proof}[Proof sketch]
The Einstein--Hilbert term gives $\Gmn/(16\pi G)$ by the standard Hilbert
variation~\cite{Wald1984}.  The matter term gives $-\Tmn/2$ from the Hilbert
prescription $\Tmn = -2\,\delta S_\Theta/(\sqrt{-g}\,\delta g^{\mu\nu})$.
Combining: $\Gmn = 8\pi G\,\Tmn$.
Symmetry of $\Tmn$: both terms are manifestly symmetric.
Conservation: follows from the Bianchi identity and Noether's theorem applied to
diffeomorphism invariance of $S_\Theta$.
See Appendix~\ref{app:tmunu} for the explicit computation.

\noindent\textit{Full proof: \texttt{canonical/t\_munu/step3\_einstein\_with\_matter.tex};
$\Tmn$ properties: \texttt{canonical/geometry/stress\_energy.tex}.}
\end{proof}

% ===========================================================================
\section{Schwarzschild Metric and ASD Condition}
\label{sec:schwarzschild}
% ===========================================================================

\subsection{The \texorpdfstring{$\Theta_0$}{Theta0} Ansatz}
\label{sec:ansatz}

\begin{definition}[Spherically symmetric vacuum ansatz]
\label{def:ansatz}
The most general spherically symmetric, time-independent admissible field in
$\Aubt$ consistent with $\Theta \to \mathbf{1}$ as $r\to\infty$ takes the form
\[
  \Theta_0 \;=\; e^{i\Phi(r)}\bigl[f(r)\,\mathbf{1} + g(r)\,\boldsymbol{e}_r\bigr],
\]
where $\boldsymbol{e}_r$ is the radial Clifford basis element and $f,g,\Phi$
are real functions of $r$.  This ansatz is the unique (up to gauge)
spherically symmetric vacuum solution of the UBT Euler--Lagrange equation; it is
\textbf{not} reverse-engineered from the known Schwarzschild form.
\end{definition}

\begin{theorem}[Schwarzschild metric recovery \PROVED]
\label{thm:schwarzschild}
With $g(r) = r\Psi(r)^2$, $f'(r) = \Psi(r)\sqrt{2M/r}$, and
$\Phi(r) = (1-M/2r)/(1+M/2r)$, the ansatz $\Theta_0$ reproduces the Schwarzschild
metric in isotropic coordinates:
\[
  g_{tt} = -\Phi(r)^2, \qquad
  g_{ij} = \Psi(r)^4\,\delta_{ij}, \qquad
  \Psi(r) = 1 + \frac{M}{2r}.
\]
The spatial sector is verified numerically to relative error $<10^{-15}$
(Appendix~\ref{app:numerics}).  The temporal component $g_{tt} = -\Phi^2$ is
recovered analytically once the full complex-time structure $\tau = t + i\psi$
is used (see Remark~\ref{rem:gtt}).
\end{theorem}

\begin{proof}
The ansatz functions must satisfy the ODE consistency condition
$f'(r)^2 + g'(r)^2 = \Psi(r)^4$.  Direct computation:
\[
  f' = \Psi\sqrt{\frac{2M}{r}}, \qquad
  g' = 1 - \frac{M^2}{4r^2}.
\]
\[
  f'^2 + g'^2
    = \Psi^2\cdot\frac{2M}{r} + \left(1-\frac{M^2}{4r^2}\right)^2
    = \Psi^4. \quad\checkmark
\]
Substituting into the metric formula \eqref{eq:real_metric} gives the stated
components.  The Schwarzschild formula is used only as a \emph{check}, not as
input (see~\texttt{known\_solution\_checks.md}, Check~1).

\noindent\textit{Full derivation: \texttt{canonical/geometry/biquaternionic\_vacuum\_solutions.tex~\S3}.}
\end{proof}

\begin{remark}[Temporal component]\label{rem:gtt}
The static real-quaternion ansatz $\Theta_0$ satisfies $\partial_t\Theta_0 = 0$,
so the metric formula \eqref{eq:real_metric} alone gives $g_{tt} = 0$.  The
Lorentzian component $g_{tt} = -\Phi^2$ is recovered via
$\partial_\psi\Theta_0 = i\Phi(r)\Theta_0$, giving
$g_{tt} = -\Re[\Tr(\partial_\psi\Theta_0\cdot\partial_\psi\Theta_0^\dagger)] = -\Phi(r)^2$,
which requires the full complex-time structure.  This is an expected feature of the
static ansatz --- AXIOM-B guarantees $g_{tt} < 0$ as a theorem (Step~3) without
numerical verification of this particular component.
\end{remark}

\subsection{Anti-Self-Dual Condition and Twistor Space}
\label{sec:asd}

\begin{theorem}[ASD condition and twistor space \PROVED]
\label{thm:asd}
For $\Theta \in \mathrm{SU}(2)_- \subset \C\otimes\Hquat$, smooth with $|\Theta|=1$:
\begin{enumerate}
  \item The holonomy of $g_{\mu\nu}[\Theta]$ lies in
    $\mathrm{Sp}(1) \cong \mathrm{SU}(2)_-$, implying the anti-self-dual (ASD)
    Weyl condition $C^+ = 0$.
  \item Combined with the vacuum UBT equation $\nabla^\dagger\nabla\Theta = 0$
    (which gives $R_{\mu\nu} = 0$ via the GR chain), the metric is ASD Ricci-flat.
  \item By the Penrose nonlinear graviton theorem~\cite{Penrose1976},
    $g_{\mu\nu}[\Theta]$ admits a curved twistor space description.
\end{enumerate}
\emph{Note}: Schwarzschild (Petrov type D) lies outside the $\mathrm{SU}(2)_-$
sector.  This is consistent: Schwarzschild has a non-zero self-dual Weyl tensor.

\noindent\textit{Source: \texttt{canonical/gr\_closure/asd\_condition\_ubt.tex}.}
\end{theorem}

\begin{remark}[Independent structural check]
Theorem~\ref{thm:asd} provides an independent structural check on the metric
derivation: the ASD Ricci-flat sector of UBT matches exactly the sector described
by Penrose's twistor construction, which is a classical result of mathematical
relativity~\cite{Penrose1976,Ward1977}.
\end{remark}

% ===========================================================================
\section{Linearised Gravity and the Regge-Wheeler Equation}
\label{sec:rw}
% ===========================================================================

\begin{theorem}[Linearised GR and Regge-Wheeler equation \PROVED]
\label{thm:rw}
Linearising the UBT field equation around a background field
$\Theta = \Theta_0 + \epsilon\,\delta\Theta$ reproduces the linearised Einstein
equations.  For odd-parity (axial) perturbations of the Schwarzschild background
decomposed into spherical harmonic modes $(\ell, m, \omega)$, the master
perturbation equation reduces to the \textbf{Regge-Wheeler equation}
\[
  \left[\frac{\mathrm{d}^2}{\mathrm{d}r_*^2} + \omega^2 - \RW(r)\right]
  \Psi_{\mathrm{RW}} \;=\; 0,
\]
where $r_* = r + 2M\ln|r/2M - 1|$ is the tortoise coordinate and
\[
  \RW(r) \;=\; \left(1-\frac{2M}{r}\right)
  \left[\frac{\ell(\ell+1)}{r^2} - \frac{6M}{r^3}\right]
\]
is the Regge-Wheeler potential for spin-2 gravitational perturbations.
No additional input beyond the metric chain (Steps~1--5) is required.
\end{theorem}

\begin{proof}[Proof sketch]
The linearised UBT field equation at first order in $\epsilon$ is
$\nabla^\dagger\nabla\,\delta\Theta = \kappa\,\delta\mathcal{T}$.  Taking the
real projection and using the chain $g \to \Gamma \to R$ applied to
$\gmn[\Theta_0 + \epsilon\,\delta\Theta]$ yields the linearised Einstein equations
$\delta G_{\mu\nu} = 8\pi G\,\delta\Tmn$.  Restricting to the odd-parity sector
and performing the standard angular decomposition~\cite{Regge1957} reduces the
system to the Regge-Wheeler equation with potential $\RW$.
\end{proof}

\begin{tcolorbox}[colback=orange!4!white,colframe=orange!55!black,
  title=Status of even-parity (Zerilli) equation]
The even-parity (polar) perturbation equation of Schwarzschild --- the Zerilli
equation~\cite{Zerilli1970} --- has \textbf{not} yet been derived from UBT.
This is GAP-Z (Section~\ref{sec:open}).  Closing this gap requires implementing
Chandrasekhar's two-potential transformation in the even-parity UBT $\Theta$
sector.  GAP-Z does \textbf{not} affect the validity of Theorem~\ref{thm:rw}
or any of the main claims.
\end{tcolorbox}

% ===========================================================================
\section{Open Problems}
\label{sec:open}
% ===========================================================================

The following open problems are precisely bounded.
\textbf{None affects the validity of Theorems~\ref{thm:metric}--\ref{thm:rw}
or the Main Theorem.}

\begin{tcolorbox}[colback=red!4!white,colframe=red!55!black,
  title=GAP-10: Off-Shell $\Theta$-Only Closure \OPEN~[L2],breakable]

\textbf{On-shell result (proved)}: For $\Theta \in \Aubt$ satisfying its own
Euler--Lagrange equation under a gauge-reduced local rank condition,
$\delta\hat{S}/\delta\Theta = 0$ is equivalent to the Einstein equations
evaluated on $g = g[\Theta]$.

\medskip
\textbf{Off-shell gap}: Global non-degeneracy of
$J = \delta g^{\mu\nu}/\delta\Theta$ for \emph{all} field configurations
(not only on-shell $\Theta \in \Aubt$) is not proved.

\medskip
\textbf{Missing lemma}: Show that $\ker J$ consists only of gauge directions
(pure phase rotation or diffeomorphism) for all $\Theta$ in the full off-shell
field space.

\medskip
\textbf{Known obstructions}:
\begin{itemize}
  \item \emph{Rank mismatch}: $\Re(\nabla^\dagger\nabla\Theta)$ is rank-0;
    $\Gmn$ is rank-2.  The multi-step chain
    $\Theta\to\partial_\mu\Theta\to\Gmn\to\gmn$ must remain non-degenerate off-shell.
  \item \emph{Topology}: Global injectivity of $\Theta\to g[\Theta]$ requires
    analysis of $H^2(M^4,\mathbb{Z})$ for the $\Theta$-bundle.
  \item \emph{Non-perturbative}: A fixed-point theorem in an appropriate Sobolev
    space is needed for off-shell well-posedness.
\end{itemize}

\medskip
\textbf{Scope}: This is a question about off-shell path-integral completeness.
The on-shell classical GR recovery is unaffected.

\medskip
\textit{Ref: \texttt{research\_tracks/T1\_GR/proof\_gap\_list.md~\S GAP-10};
\texttt{canonical/gr\_closure/step2\_theta\_only\_closure.tex}.}
\end{tcolorbox}

\begin{tcolorbox}[colback=orange!4!white,colframe=orange!55!black,
  title=GAP-Z: Zerilli Equation (Even-Parity Graviton) \OPEN~[L2]]

\textbf{Proved}: The Regge-Wheeler equation (odd-parity graviton) is derived
(Theorem~\ref{thm:rw}).

\medskip
\textbf{Missing}: The Zerilli equation for even-parity (polar) perturbations:
\[
  \left[\frac{\mathrm{d}^2}{\mathrm{d}r_*^2} + \omega^2 - V_{\mathrm{Z}}(r)\right]
  \Psi_{\mathrm{Z}} = 0.
\]
Even-parity modes couple scalar and tensor sectors in a way requiring
Chandrasekhar's two-potential transformation~\cite{Chandrasekhar1983}, which has
not yet been implemented for the UBT even-parity $\Theta$ sector.

\medskip
\textbf{Path to closure}: (i) Derive the even-parity linearised UBT field
equation; (ii) show it reduces to Zerilli via Chandrasekhar's transformation.
This is the highest-priority open item in the graviton sector.

\medskip
\textit{Ref: \texttt{research\_tracks/T1\_GR/proof\_gap\_list.md~\S GAP-Z}.}
\end{tcolorbox}

\begin{tcolorbox}[colback=gray!5!white,colframe=gray!45!black,
  title=Additional Lower-Priority Gaps (for completeness)]
\begin{itemize}
  \item \textbf{GAP-C}: Cosmological (FRW/de~Sitter) solutions not yet
    derived from an explicit time-dependent $\Theta$ ansatz.
  \item \textbf{GAP-M}: Compact $M^4$ off-shell closure requires bundle
    topology analysis.
  \item \textbf{GAP-Q}: Path-integral quantisation of UBT --- long-term,
    depends on GAP-10.
\end{itemize}
None of these block the on-shell classical result.
\end{tcolorbox}

% ===========================================================================
\section{Discussion}
\label{sec:discussion}
% ===========================================================================

\subsection{Summary}
\label{sec:summary}

The chain
\[
  \Theta \;\longrightarrow\; \gmn \;\longrightarrow\;
  \Gamma^\lambda_{\mu\nu} \;\longrightarrow\;
  R^\rho{}_{\sigma\mu\nu} \;\longrightarrow\;
  \Gmn = 8\pi G\,\Tmn
\]
is closed at the [L1] level.  UBT contains standard GR as an exact on-shell sector:
the metric and Lorentzian signature are derived quantities, not postulates; the
Schwarzschild metric is reproduced to numerical precision; the Regge-Wheeler
equation follows without additional input.  Two explicitly bounded open problems
(GAP-10 and GAP-Z) do not affect this result.

\subsection{Comparison with Existing Algebraic Approaches to GR}
\label{sec:comparison}

\begin{center}
\begin{tabular}{p{3.5cm}p{3.5cm}p{3.5cm}p{3.5cm}}
\toprule
\textbf{Feature} & \textbf{UBT (this paper)} &
\textbf{Prior biquaternion gravity}~\cite{Adler1995,DeLeo1996} &
\textbf{Spectral action}~\cite{ConnesLott1990} \\
\midrule
Metric derivation &
  Derived from $\Theta$ via bilinear formula &
  Postulated or imposed &
  Derived from Dirac operator \\
Lorentzian signature &
  Proved (Thm.~\ref{thm:signature}) &
  Assumed &
  Assumed \\
Free parameters in GR chain &
  None &
  Typically free couplings &
  None (but different algebra) \\
Algebra dimension &
  $\dim_\R\B = 8$ &
  Varies &
  $\dim_\R A = 21$ \\
Schwarzschild recovery &
  Analytical + numerical &
  Not demonstrated &
  Not demonstrated \\
Regge-Wheeler equation &
  Proved &
  Not addressed &
  Not addressed \\
\bottomrule
\end{tabular}
\end{center}

The key structural distinction: in UBT, $\gmn$ is the pullback of the Clifford
norm by $\Theta$, which is analogous to the Penrose sigma-model
$\gmn = \Re[\partial_\mu\Theta\cdot\partial_\nu\Theta^\dagger]$.  No prior
biquaternion gravity paper has this property.

\subsection{Relation to GR Predictions}
\label{sec:gr_predictions}

Since UBT embeds GR exactly in the real-sector limit, every experimental
confirmation of GR --- perihelion precession, gravitational lensing, frame
dragging, gravitational waves detected by LIGO/Virgo~\cite{Abbott2016} ---
is automatically a confirmation of the UBT real sector.  The biquaternionic
imaginary sector introduces additional degrees of freedom (phase curvature) that
are invisible to classical measurements because ordinary matter couples only to
the real metric $\gmn$.

\subsection{Outlook}
\label{sec:outlook}

The following future directions follow directly from this work:

\begin{enumerate}
  \item \textbf{GAP-Z (Zerilli equation)}: Complete the graviton sector by
    deriving the even-parity equation.  Clear path via Chandrasekhar's
    transformation.

  \item \textbf{GAP-C (Cosmological solutions)}: Derive the Friedmann equations
    from a time-dependent spherically symmetric $\Theta$ ansatz.  The procedure
    is the same as for Schwarzschild with $\partial_t\Theta \neq 0$.

  \item \textbf{T2\_GAUGE (Gauge sector)}: The companion paper derives the
    $\mathrm{SU}(3)\times\mathrm{SU}(2)_L\times\mathrm{U}(1)_Y$ gauge structure
    from the imaginary sector of $\B$.

  \item \textbf{GAP-10 (Off-shell closure)}: Long-term topological problem.
    Approach: compute cohomology of $\ker J$ sheaf over $M^4$.
\end{enumerate}

% ===========================================================================
%  PROOF STATUS TABLE
% ===========================================================================

\begin{tcolorbox}[colback=blue!3!white,colframe=blue!55!black,
  title=T1\_GR Proof Status Summary]
\begin{center}
\begin{tabular}{clll}
\toprule
\textbf{Step} & \textbf{Claim} & \textbf{Status} & \textbf{Source} \\
\midrule
1  & Metric $\gmn$ from $\Theta$               & \PROVED & step1\_metric\_bridge.tex \\
2  & Non-degeneracy $\det(g)\neq 0$            & \PROVED & step2\_nondegeneracy.tex \\
3  & Signature $(-,+,+,+)$ from AXIOM-B       & \PROVED & step3\_signature\_theorem.tex \\
4  & $g\to\Gamma\to R$ geometric chain         & Standard GR & curvature.tex \\
5  & Einstein eqs $\Gmn=8\pi G\Tmn$           & \PROVED & step3\_einstein\_with\_matter.tex \\
6a & Schwarzschild metric (spatial)            & \PROVED & verify\_schwarzschild\_theta.py \\
6b & ASD condition / twistor                  & \PROVED & asd\_condition\_ubt.tex \\
6c & Regge-Wheeler equation                   & \PROVED & (linearised GR chain) \\
7a & Zerilli equation (even-parity)            & \OPEN~[L2] & — \\
7b & Off-shell $\Theta$-only closure (GAP-10) & \OPEN~[L2] & step2\_theta\_only\_closure.tex \\
\bottomrule
\end{tabular}
\end{center}
\end{tcolorbox}

% ===========================================================================
%  APPENDICES
% ===========================================================================

\appendix

% ===========================================================================
\section{Full Proof of the Lorentzian Signature Theorem}
\label{app:signature}
% ===========================================================================

This appendix gives the complete algebraic argument for
Theorem~\ref{thm:signature}.

\textbf{Step A1: Clifford identification.}
The biquaternion algebra satisfies
$\B = \C\otimes_\R\Hquat \cong \mathrm{Cl}_{1,3}(\R)$ as real algebras.
Under this isomorphism the generators split into one timelike generator $\gamma^0$
and three spacelike generators $\gamma^i$ ($i=1,2,3$) satisfying
\[
  (\gamma^0)^2 = -1, \qquad (\gamma^i)^2 = +1, \qquad
  \{\gamma^\mu,\gamma^\nu\} = 2\eta^{\mu\nu}.
\]
In the matrix representation ($\B \cong \mathrm{Mat}(2,\C)$):
$\gamma^0 = i\sigma^0 = iI_2$, $\gamma^i = \sigma^i$ (Pauli matrices).

\textbf{Step A2: AXIOM-B constraint.}
AXIOM-B (Assumption~A2) states $\langle\pT,\pT\rangle_\eta < 0$.
Under $\tau = t + i\psi$:
$\pT = \partial_t + i\partial_\psi$.
Since $\pT$ lies in the timelike Clifford sector,
$\pT^2 < 0$ in the Clifford quadratic form $q(\gamma^0)^2 + (\gamma^i)^2 = -1+1+1+1$.

\textbf{Step A3: Temporal metric component.}
For $\partial_0 = \partial_t$ in the timelike $\gamma^0$-sector, writing
$\partial_0\Theta = \xi\,\gamma^0\Theta_0$ for some scalar $\xi$:
\[
  \Re\bigl[\Tr(\partial_0\Theta\cdot\partial_0\Theta^\dagger)\bigr]
  = |\xi|^2\,\Re\bigl[\Tr((\gamma^0)^2)\bigr]
  = |\xi|^2\,\Re[\Tr(-I_2)]
  = -2|\xi|^2 < 0.
\]
Setting $\mathcal{N} = -\Re[\Tr(\partial_0\Theta\cdot\partial_0\Theta^\dagger)]
= 2|\xi|^2 > 0$ gives $g_{00} = -1 < 0$.

\textbf{Step A4: Spatial metric components.}
For $\partial_i\Theta = \zeta\,\gamma^i\Theta_0$:
\[
  \Re\bigl[\Tr(\partial_i\Theta\cdot\partial_i\Theta^\dagger)\bigr]
  = |\zeta|^2\,\Re[\Tr((\gamma^i)^2)]
  = |\zeta|^2\,\Re[\Tr(I_2)]
  = 2|\zeta|^2 > 0,
\]
giving $g_{ii} > 0$.

\textbf{Step A5: Off-diagonal components.}
Mixed products $\gamma^0(\gamma^i)^\dagger + \gamma^i(\gamma^0)^\dagger$
are traceless in the $\mathrm{Mat}(2,\C)$ representation, so $g_{0i} = 0$ in
the static spatially isotropic sector.  General rotating spacetimes may have
$g_{0i} \neq 0$ (frame dragging), consistent with the full Lorentzian block
structure $(-;+,+,+)$.

\textbf{Conclusion.}
The eigenvalue structure of $\gmn$ has exactly one negative eigenvalue
($g_{00} < 0$) and three positive eigenvalues ($g_{ii} > 0$), i.e.\ signature
$(-,+,+,+)$.  This follows entirely from AXIOM-B without any assumption about
the specific form of $\Theta$.

\medskip
\begin{center}
\begin{tabular}{lll}
\toprule
\textbf{Framework} & \textbf{Signature input} & \textbf{Status} \\
\midrule
Standard GR        & Lorentzian signature assumed directly & Postulate \\
String theory      & Target-space metric assumed Lorentzian & Postulate \\
Loop Quantum Gravity & Signature encoded in spin-foam amplitudes & Structural input \\
UBT (this paper)   & AXIOM-B $\Rightarrow$ signature is a theorem & Derived \\
\bottomrule
\end{tabular}
\end{center}
\captionof{table}{Signature input in competing frameworks.}
\label{tab:signature_comparison}

\medskip
\textit{Full algebraic proof with explicit $2\times 2$ matrix computations:
\texttt{canonical/gr\_closure/step3\_signature\_theorem.tex}, Theorem~2.}

% ===========================================================================
\section{Stress-Energy Tensor and Conservation}
\label{app:tmunu}
% ===========================================================================

\subsection{Derivation of $T_{\mu\nu}$}

The $\Theta$-matter action with kinetic term is
\[
  S_\Theta[g,\Theta] \;=\; \int\!\sqrt{-g}\;
  \Re\bigl[\Tr((D_\mu\Theta)^\dagger D^\mu\Theta)\bigr]\,\mathrm{d}^4x.
\]
The Hilbert prescription $\Tmn = -2\,\delta S_\Theta/(\sqrt{-g}\,\delta g^{\mu\nu})$
gives
\[
  \Tmn \;=\; \Re\bigl[\Tr(\partial_\mu\Theta\,\partial_\nu\Theta^\dagger)\bigr]
  - \tfrac{1}{2}g_{\mu\nu}\,g^{\alpha\beta}
    \Re\bigl[\Tr(\partial_\alpha\Theta\,\partial_\beta\Theta^\dagger)\bigr].
\]

\subsection{Symmetry}

Both terms in $\Tmn$ are manifestly symmetric: the trace bilinear is symmetric
by the same argument as in Theorem~\ref{thm:metric}, and $g_{\mu\nu}$ is
symmetric by Theorem~\ref{thm:metric}.  Hence $\Tmn = T_{\nu\mu}$.

\subsection{Conservation}

\textbf{Method 1 (Bianchi identity):}
The Bianchi identity $\nabla^\mu\Gmn = 0$ combined with the Einstein equations
$\Gmn = 8\pi G\Tmn$ immediately gives $\nabla^\mu\Tmn = 0$.

\textbf{Method 2 (Noether's theorem):}
The action $S_\Theta$ is invariant under diffeomorphisms $x\mapsto x + \xi(x)$.
By Noether's second theorem, this invariance implies $\nabla^\mu\Tmn = 0$
identically, independently of the field equation.

\textit{Full proof: \texttt{canonical/geometry/stress\_energy.tex}.}

% ===========================================================================
\section{Numerical Verification of the Schwarzschild Metric}
\label{app:numerics}
% ===========================================================================

\subsection{Setup}

Script: \texttt{tools/verify\_schwarzschild\_theta.py}.
Mass parameter: $M = 1$ (geometrised units).
Ansatz functions evaluated at each $r$; metric components computed from
formula~\eqref{eq:real_metric}; results compared with the isotropic Schwarzschild
metric $\Psi(r)^4\,\delta_{ij}$.

\subsection{ODE Consistency Condition}

\begin{center}
\begin{tabular}{rrrrrr}
\toprule
$r/M$ & $f'$ & $g'$ & $f'^2+g'^2$ & $\Psi^4$ & Error \\
\midrule
2.0   & 1.250000 & 0.937500 & 2.441406 & 2.441406 & $0.00$ \\
5.0   & 0.695701 & 0.990000 & 1.464100 & 1.464100 & $1.5\times10^{-16}$ \\
10.0  & 0.469574 & 0.997500 & 1.215506 & 1.215506 & $1.8\times10^{-16}$ \\
100.0 & 0.142128 & 0.999975 & 1.020151 & 1.020151 & $4.3\times10^{-16}$ \\
\bottomrule
\end{tabular}
\end{center}

\subsection{Spatial Metric Components}

\begin{center}
\begin{tabular}{rrrrrrr}
\toprule
$r/M$ & $\Psi^4$ (expected) & $g_{xx}$ & $g_{yy}$ & $g_{zz}$ & $\max|g_\mathrm{off}|$ & Status \\
\midrule
2.0  & 2.441406 & 2.441406 & 2.441406 & 2.441406 & $5.6\times10^{-17}$ & \checkmark \\
5.0  & 1.464100 & 1.464100 & 1.464100 & 1.464100 & $1.9\times10^{-16}$ & \checkmark \\
10.0 & 1.215506 & 1.215506 & 1.215506 & 1.215506 & $2.8\times10^{-17}$ & \checkmark \\
50.0 & 1.040604 & 1.040604 & 1.040604 & 1.040604 & $7.3\times10^{-17}$ & \checkmark \\
100.0& 1.020151 & 1.020151 & 1.020151 & 1.020151 & $2.7\times10^{-16}$ & \checkmark \\
\bottomrule
\end{tabular}
\end{center}

\textbf{Result}: All spatial components agree with $\Psi^4\,\delta_{ij}$ to
floating-point precision ($\sim10^{-15}$ relative error).
Off-diagonal components are numerically zero.

\subsection{Honest Accounting}

\begin{center}
\begin{tabular}{llll}
\toprule
Component & Analytical & Numerical & Status \\
\midrule
Metric formula (Steps 1--3) & \checkmark & -- & Proved \\
Levi-Civita + curvature (Step 4) & \checkmark & -- & Standard GR \\
Einstein equations (Step 5) & \checkmark & -- & Proved \\
Schwarzschild $g_{ij}$ (spatial) & \checkmark (exact) & \checkmark ($<10^{-15}$) & Verified \\
Schwarzschild $g_{i0}$ (off-diag) & \checkmark (symmetry) & \checkmark (zero) & Verified \\
Schwarzschild $g_{tt}$ (temporal) & \checkmark (complex $\tau$) & -- & Analytically known \\
ODE consistency $f'^2+g'^2=\Psi^4$ & \checkmark (algebraic) & \checkmark (FP) & Verified \\
ASD Weyl condition & \checkmark & -- & Proved \\
Twistor space & \checkmark (Penrose) & -- & Proved \\
Regge-Wheeler (odd-parity) & \checkmark & -- & Proved \\
Zerilli (even-parity) & -- & -- & Open [GAP-Z] \\
\bottomrule
\end{tabular}
\end{center}

\subsection{Reproduction}

\begin{verbatim}
pip install numpy
python tools/verify_schwarzschild_theta.py
# optional: python tools/verify_schwarzschild_theta.py --mass 2.0 --r_values 3,6,12
\end{verbatim}

Exit code 0: all spatial components agree within tolerance $10^{-8}$ (default).

% ===========================================================================
%  BIBLIOGRAPHY
% ===========================================================================

\begin{thebibliography}{99}

\bibitem{Wald1984}
R.~M.~Wald,
\textit{General Relativity}.
University of Chicago Press, 1984.

\bibitem{MTW1973}
C.~W.~Misner, K.~S.~Thorne, J.~A.~Wheeler,
\textit{Gravitation}.
W.~H.~Freeman, 1973.

\bibitem{Adler1995}
S.~L.~Adler,
\textit{Quaternionic Quantum Mechanics and Quantum Fields}.
Oxford University Press, 1995.

\bibitem{DeLeo1996}
S.~De Leo,
``Quaternions and special relativity,''
\textit{J. Math. Phys.}\ \textbf{37}, 2955 (1996).

\bibitem{Finkelstein1962}
D.~Finkelstein, J.~M.~Jauch, S.~Schiminovich, D.~Speiser,
``Foundations of quaternion quantum mechanics,''
\textit{J. Math. Phys.}\ \textbf{3}, 207 (1962).

\bibitem{Penrose1976}
R.~Penrose,
``Non-linear gravitons and curved twistor theory,''
\textit{Gen. Rel. Grav.}\ \textbf{7}, 31 (1976).

\bibitem{Ward1977}
R.~S.~Ward,
``On self-dual gauge fields,''
\textit{Phys. Lett.}\ \textbf{A61}, 81 (1977).

\bibitem{ConnesLott1990}
A.~Connes, J.~Lott,
``Particle models and noncommutative geometry,''
\textit{Nucl. Phys. Proc. Suppl.}\ \textbf{18B}, 29 (1990).

\bibitem{Regge1957}
T.~Regge, J.~A.~Wheeler,
``Stability of a Schwarzschild singularity,''
\textit{Phys. Rev.}\ \textbf{108}, 1063 (1957).

\bibitem{Zerilli1970}
F.~J.~Zerilli,
``Effective potential for even-parity Regge-Wheeler gravitational perturbation equations,''
\textit{Phys. Rev. Lett.}\ \textbf{24}, 737 (1970).

\bibitem{Chandrasekhar1983}
S.~Chandrasekhar,
\textit{The Mathematical Theory of Black Holes}.
Oxford University Press, 1983.

\bibitem{HornJohnson2013}
R.~A.~Horn, C.~R.~Johnson,
\textit{Matrix Analysis}, 2nd ed.
Cambridge University Press, 2013.

\bibitem{Abbott2016}
B.~P.~Abbott et al.\ (LIGO Scientific and Virgo Collaborations),
``Observation of gravitational waves from a binary black hole merger,''
\textit{Phys. Rev. Lett.}\ \textbf{116}, 061102 (2016).

\end{thebibliography}

\end{document}
